\documentclass[11pt]{article}

\usepackage[margin=1in]{geometry}
\usepackage{amsmath,amssymb,amsthm,mathtools}
\usepackage{enumitem}
\usepackage{booktabs}
\usepackage{longtable}
\usepackage{array}
\usepackage{xcolor}
\usepackage[hidelinks]{hyperref}
\usepackage{microtype}

\newcommand{\ALCIRCC}{\mathcal{ALCI}_{\mathrm{RCC5}}}
\newcommand{\PP}{\mathsf{PP}}
\newcommand{\PPI}{\mathsf{PPI}}
\newcommand{\PO}{\mathsf{PO}}
\newcommand{\DR}{\mathsf{DR}}
\newcommand{\EQ}{\mathsf{EQ}}
\newcommand{\Base}{\{\DR,\PO,\PP,\PPI,\EQ\}}
\newcommand{\Cl}{\operatorname{Cl}}
\newcommand{\Tp}{\operatorname{Tp}}
\newcommand{\Need}{\operatorname{Need}}
\newcommand{\reach}{\rightsquigarrow}
\newcommand{\parQ}{\mathsf{par}_Q}
\newcommand{\eqcong}{\equiv_{\EQ}}
\newcommand{\blk}{\equiv_{\mathrm{blk}}}
\newcommand{\Mate}{\operatorname{Mate}}
\newcommand{\st}{\operatorname{st}}
\newcommand{\comp}{\operatorname{comp}}
\newcommand{\orig}{\operatorname{orig}}

\newtheorem{obligation}{Proof obligation}
\newtheorem{recommendation}{Recommended repair}

\title{Formal response to the v2 \texorpdfstring{$\ALCIRCC$}{ALCI-RCC5} updates}
\author{Review memorandum}
\date{May 2026}

\begin{document}
\maketitle

\begin{abstract}
This note reviews the latest v2 documents, in particular
the completeness-extraction v2 paper and the sibling-interface v2 paper, with brief comments on the overview paper.  The main conclusion is positive but conditional.
The v2 revision correctly removes the earlier density/Hasse-cover problem by
introducing witness-generated submodels and generated cover edges.  It also
moves the proof much closer to the intended split-forest architecture.  However,
several proof obligations remain open before the decidability theorem should be
regarded as established.  The main remaining issues concern the separation of
blocking from semantic equality, the treatment of multiple incomparable upward
\(\PP\)-witnesses, explicit handling of \(\EQ\)-modalities, finite mosaic arity,
placement of \(\DR/\PO\) witnesses, and simultaneous realization of transitive
eventualities.
\end{abstract}

\section{Scope and verdict}

This response uses line numbers from the supplied TeX files where useful.  The main references are the completeness-extraction v2 source and the sibling-interface v2 source.

The v2 update is a substantial improvement over the earlier draft.  In
particular, the problematic reading of the split tree as a Hasse diagram of an
arbitrary ambient \(\PP/\PPI\)-order has been replaced by the correct generated
model picture.  The witness-generated submodel lemma, generated cover edges,
self-loops as blocking cycles, and support-closed mosaics are all steps in the
right direction.

Nevertheless, I would not yet state the main theorem as a completed proof.
The current version should be described as a strong proof programme with several
remaining formal proof obligations.  The most serious of these is the interaction
between weak-\(\EQ\) splitting, finite blocking, and multiple upward
\(\PP\)-witness requirements.

\section{What is now fixed}

\subsection{The density/Hasse issue is essentially resolved}

The completeness extraction now starts with witness-generated submodels.  The
closure definition is closed under NNF complements
(completeness v2, lines 228--235), and the
witness-generated submodel lemma is stated and proved at lines 314--351.  The
proof correctly handles the converse direction for universals by retaining a
witness for \(\exists R.\overline F\) when \(\forall R.F\) fails; see lines
344--350.

This is the right bridge between ordinary abstract satisfiability and the sparse
model class originally intended for the split-forest construction.  It shows
that arbitrary dense or non-Hasse features of a starting semantic frame are not
relevant to satisfiability of the closure formulas.  One may restrict to a
countable submodel generated by selected witnesses.

The document also now distinguishes generated cover edges from semantic Hasse
covers.  Lines 405--419 explicitly say that the generated cover relation is not
assumed to be a Hasse relation in the original ambient order.  This is the
right conceptual correction.  The proof no longer needs to read off a transitive
reduction of an arbitrary \(\PP\)-order.

\subsection{The ambient-root assumption has been removed}

The sibling-interface v2 document explicitly replaces the old ambient-root
requirement.  It says that v1 Definition 1.5 is reinterpreted so that there is
no ambient root, generated cover edges are selected witness covers, and parent
self-loops are permitted; see
sibling-interface v2, lines 167--178 and 257--289.

Likewise, the completeness extraction allows self-loops in the parent function:
\[
  \parQ(\sigma)=\sigma .
\]
This is described at lines 592--607 of the completeness file.  The intended
meaning is correct: a quotient self-loop represents infinitely many fresh
unfolding occurrences with the same finite profile.  It is not a semantic cycle
and not an \(\EQ\)-identification.

This handles the earlier stress formula
\[
  C_\uparrow = \exists\PP.\top \sqcap \forall\PP.\exists\PP.\top
\]
under the containment orientation: each occurrence has a strict later
superpart/ancestor, while the finite quotient uses a single blocking state.

\subsection{The raw pair-table problem is addressed in the right direction}

The v2 extraction replaces raw pair tables by support-closed mosaics; see lines
616--720 of the completeness file.  This is the correct response to the witness
conflation problem in the old C5-style reasoning.  A mosaic records a joint
atomic network on one finite vertex set, rather than aggregating independent
pairwise relations that may not be jointly realized.

This is a real improvement.  It preserves the original sibling-interface idea
while strengthening it enough to have a plausible completeness theorem.

\section{Remaining blockers}

\subsection{Blocking and semantic \texorpdfstring{$\EQ$}{EQ} are still not cleanly separated}

This is the most important remaining issue.

The profile definition records a typed-\(\EQ\) colour as a finite component of a
rank profile; see lines 525--545.  The finite-index lemma then says these colours
are drawn from a finite palette; see lines 551--566.  But the extracted
congruence proof says that the colour is computed from the origin, choosing one
colour per \(\EQ\)-class in the witness-generated model; see lines 864--867.
Under strong \(\EQ\), every distinct semantic element is its own \(\EQ\)-class.
A witness-generated model may be countably infinite, so one colour per semantic
\(\EQ\)-class is not a finite profile component.

There is also a direct conflict with blocking loops.  Lines 592--607 correctly
state that a quotient self-loop unfolds to fresh semantic occurrences.  Lines
845--855 also say that repeated occurrences of the same profile are not
\(\EQ\)-related.  However, the mosaic definition says that endpoints with the
same typed-\(\EQ\) colour have label \(\EQ\); see lines 693--694.  If the same
finite profile and colour are reused on successive laps of a blocking loop, this
would incorrectly collapse the infinite \(\PP\)-chain into semantic equality.

\begin{obligation}[Separate blocking from equality]
The proof must define two distinct relations:
\[
  \blk \quad\text{finite blocking/profile repetition,}
\]
\[
  \eqcong \quad\text{semantic equality / weak-}\EQ\text{ mate relation.}
\]
Equality of finite profiles, or equality of a reusable profile colour, must
never by itself imply \(\eqcong\).
\end{obligation}

\begin{recommendation}[Revise typed-\(\EQ\) colours]
Do not say that a typed-\(\EQ\) colour is one colour per semantic \(\EQ\)-class
of \(\mathcal I_w\).  Instead, make equality data explicit and local to the
finite certificate.  For example:
\begin{itemize}[leftmargin=2em]
\item The quotient contains an explicit finite relation \(\equiv_{\EQ,Q}\) on
  named equality ports, core copies, or mate slots.
\item Blocking states may repeat the same profile, but every unfolding lap gets
  fresh equality ports unless the certificate explicitly links them by
  \(\equiv_{\EQ,Q}\).
\item The mosaic condition should be changed from ``same colour implies
  \(\EQ\)'' to ``\(\EQ\)-label iff the endpoints are identified by the explicit
  equality relation visible in that mosaic.''
\item Add a prohibition that no two occurrences connected by a strict vertical
  \(\PP/\PPI\)-path may be \(\EQ\)-identified.
\end{itemize}
\end{recommendation}

With this change, blocking is purely a finite representation device, while
strong \(\EQ\) is handled only by the explicit typed congruence.

\subsection{Multiple upward \texorpdfstring{$\PP$}{PP} witnesses need an EQ-aware construction}

The containment orientation uses parent edges for superparts.  Therefore
\(\exists\PP.D\) requirements point upward.  A single occurrence in a tree has
only one parent chain.  But a semantic element may need several incomparable
proper-superpart witnesses.

A useful stress formula is
\[
\begin{aligned}
C_{AB} := {}& \exists\PP.A \sqcap \exists\PP.B \\
&\sqcap\ \forall\PP\bigl(\neg A \sqcup
   (\forall\PP.\neg B \sqcap \forall\PPI.\neg B)\bigr) \\
&\sqcap\ \forall\PP\bigl(\neg B \sqcup
   (\forall\PP.\neg A \sqcap \forall\PPI.\neg A)\bigr).
\end{aligned}
\]
This concept is satisfiable by an element \(x\) with two incomparable proper
superparts \(y\) and \(z\), where \(y\models A\), \(z\models B\), and
\(y\PO z\).  The universal conjuncts prevent the \(A\)- and \(B\)-witnesses from
being comparable along one parent chain.

Thus a single occurrence of \(x\) cannot realize both \(\PP\)-witnesses by one
linear parent chain.  The original split-copy idea is designed exactly for this
case: create weak-\(\EQ\) copies of \(x\), put one copy under the \(A\)-superpart
and another copy under the \(B\)-superpart, and then recover strong \(\EQ\) by a
typed congruence.

The current v2 reachability clause is still too local.  It says, essentially,
\[
  \exists R.D\in\lambda_Q(\sigma)
  \quad\Longrightarrow\quad
  \exists\sigma'\; (\sigma\reach_R\sigma'\text{ and }D\in\lambda_Q(\sigma'))
\]
for \(R\in\{\PP,\PPI\}\); see lines 960--979 and 1067--1075.  This does not
say that a witness may be realized along the parent chain of an \(\EQ\)-mate of
\(\sigma\).

\begin{obligation}[EQ-aware reachability]
For split copies, existential and universal \(\PP/\PPI\) clauses must be checked
on equality classes, not on a single occurrence profile.  A suitable
existential condition is:
\[
  \exists\PP.D\in\lambda(\sigma)
  \Rightarrow
  \exists\widehat\sigma\in\Mate(\sigma)\;\exists\sigma'
  \bigl(\widehat\sigma\reach_{\PP}\sigma'\wedge D\in\lambda(\sigma')\bigr),
\]
and similarly for \(\PPI\).  A suitable universal condition is:
\[
  \forall\PP.D\in\lambda(\sigma)
  \Rightarrow
  \forall\widehat\sigma\in\Mate(\sigma)\;\forall\sigma'
  \bigl(\widehat\sigma\reach_{\PP}\sigma'\Rightarrow D\in\lambda(\sigma')\bigr).
\]
Here \(\Mate(\sigma)\) must be derived from explicit typed-\(\EQ\) mate data,
not from profile equality.
\end{obligation}

This repair is needed for completeness.  After quotienting the weak split model
by \(\eqcong\), a witness found above any mate copy becomes a witness for the
strong semantic element.

\subsection{The role \texorpdfstring{$\EQ$}{EQ} must be normalized or handled explicitly}

The language treats the five RCC5 atoms as modal roles, including \(\EQ\).  The
validity proof distinguishes only \(\DR/\PO\) witnesses and \(\PP/\PPI\) witnesses;
see lines 1067--1075.  There is no case for \(R=\EQ\).

Under strong equality semantics this is easy to fix.  Since
\[
  \exists\EQ.D \equiv D,
  \qquad
  \forall\EQ.D \equiv D,
\]
all \(\EQ\)-modalities may be eliminated during preprocessing, or the type
conditions may explicitly enforce these equivalences.

\begin{recommendation}[EQ-normalization]
Add a preprocessing lemma: for every concept \(C\), compute in polynomial time a
concept \(C^{-\EQ}\) in which \(\exists\EQ.D\) and \(\forall\EQ.D\) are replaced
by \(D\), recursively, and prove equivalence under strong \(\EQ\).  Then the
rest of the proof may assume existential roles are only
\(\DR,\PO,\PP,\PPI\).
\end{recommendation}

Alternatively, add \(\EQ\) as a third case in the validity proof.

\subsection{The finite mosaic arity bound is not yet proved}

The support-closed mosaic idea is sound in spirit, but the finite arity bound is
not yet formal.  Lines 722--733 say mosaics may be bounded by
\(|\Cl(C_0)|+4\) vertices because larger mosaics decompose into overlapping
smaller ones.  Lines 1258--1264 then acknowledge that a rigorous treatment would
derive this bound from a Ramsey-type argument.

This is not just a cosmetic matter.  The decidability proof depends on an
effectively enumerable finite certificate class.  Therefore the arity bound must
be a theorem, not an informal statement.

\begin{obligation}[Finite mosaic locality]
Prove that there is an effectively computable number \(N(C_0)\) such that every
finite prefix network needed by the split-forest soundness construction is
controlled by support-closed mosaics of arity at most \(N(C_0)\).  Equivalently,
show that all required restriction, triple-composition, typed-\(\EQ\)
replacement, containment-extension, and witness-support checks are local below
that bound.
\end{obligation}

One possible route is to define \(N(C_0)\) from the number of endpoint summaries
and closure-existential requirements, rather than asserting \(|\Cl(C_0)|+4\).
A second route is to state and prove a patchwork/locality theorem for abstract
RCC5 mosaics with typed endpoint summaries.

\subsection{Placement of \texorpdfstring{$\DR/\PO$}{DR/PO} witnesses remains schematic}

The witness-menu construction for \(\DR/\PO\) existentials is at lines 775--795.
It says that the selected witness occurrence exists in the split tree, possibly
as a sibling under a common ancestor.  But the construction that places selected
\(\DR/\PO\) witnesses into side contexts or sibling mosaics is not given.

For completeness, the proof needs to explain how every selected \(\DR\)- or
\(\PO\)-witness from the witness-generated submodel is represented in the
containment split presentation while preserving:
\begin{itemize}[leftmargin=2em]
\item the chosen \(\DR/\PO\) relation to the source occurrence;
\item all universal restrictions at the source and witness;
\item all composition constraints with the vertical ancestors and descendants;
\item consistency with sibling mosaics and typed-\(\EQ\) mate data.
\end{itemize}

\begin{recommendation}[Side-context witness lemma]
Add a lemma of the following form.

\medskip
\noindent\emph{Lemma.}  Let \(u\) be an occurrence and let
\(b=w(\orig(u),\exists R.D)\) with \(R\in\{\DR,\PO\}\).  Then the generated split
presentation can be extended, by adding a bounded side context and a
support-closed mosaic, with an occurrence \(v\) of origin \(b\) such that
\(\rho(u,v)=R\), \(D\in\lambda(v)\), and all relation-safety and RCC5 composition
constraints involving the finite interface are preserved.
\end{recommendation}

Without such a lemma, the witness-menu entries are not yet backed by a formal
model-to-certificate construction.

\subsection{Reachability discharge must be simultaneous and coherent}

The reachability formulation is much cleaner than the earlier automata-heavy
presentation.  In a deterministic parent graph, a \(\PP\)-eventuality can often
be reduced to reachability along the unique parent walk.  However, the paper
itself notes that simultaneous discharge requires a single coherent infinite
walk; see lines 1266--1274.

The current argument says this follows from finiteness and a patchwork property.
That needs to be a proof, especially once split copies and \(\EQ\)-aware
reachability are introduced.  The issue is not just whether each eventuality has
some reachable witness profile, but whether all eventualities required along one
unfolded branch can be realized together in the same unfolding.

\begin{obligation}[Coherent eventuality realization]
Add a finite-cycle realization lemma.  A possible statement is:

\medskip
\noindent\emph{Lemma.}  Let \(G_Q\) be the finite parent graph equipped with
explicit equality-mate data.  If every \(\PP/\PPI\)-eventuality satisfies the
EQ-aware reachability condition and every universal is closed under the relevant
reachability relation, then the canonical unfolding of \(G_Q\) realizes all
vertical eventualities simultaneously.

\medskip
For \(\PP\) under a functional parent map, this can be proved by analyzing the
unique ultimately periodic parent walk.  For \(\PPI\), where inverse-parent
branching may be non-deterministic, the proof must use the child multiplicity
component of the profile to show that each occurrence has the required child
branch witnesses.
\end{obligation}

\subsection{Minor source issue: undefined \texttt{\\sst}}

The completeness source uses \verb|\sst| in the status partition at lines
632--638, but no macro \verb|\sst| appears to be defined.  This causes an
undefined-control-sequence error under ordinary LaTeX compilation.  The source
should either define
\[
  \verb|\newcommand{\sst}{\operatorname{st}}|
\]
or replace \verb|\sst| with the already defined status notation.

\section{Assessment of the companion sibling-interface v2 paper}

The sibling-interface v2 note is useful and should be kept as a delta document.
It correctly states that the old soundness chain is retained while the ambient
root and v1 decidability sketch are replaced.  In particular, the explicit list
of retained and replaced material at lines 131--180 is helpful.

However, the sibling-interface note inherits the same unresolved obligations
from the completeness paper.  In particular:
\begin{enumerate}[leftmargin=2em]
\item Self-loops must be blocking only, never equality.  The note says this at
  lines 409--414, but the full finite certificate still needs explicit equality
  data that cannot identify blocking laps.
\item Reachability discharge at lines 459--470 must be made EQ-aware to handle
  split copies with several incomparable upward \(\PP\)-witnesses.
\item The theorem at lines 534--580 relies on the companion completeness theorem;
  therefore it should be stated conditionally until the obligations in this
  response are closed.
\end{enumerate}

A good wording would be:
\begin{quote}
Together with a complete proof of the v2 completeness-extraction theorem, the
preserved v1 soundness chain yields decidability.
\end{quote}
Rather than:
\begin{quote}
Combined with v2 completeness extraction, this establishes decidability.
\end{quote}
The latter is too strong until the completeness extraction is fully formalized.

\section{Comments on the overview document}

The overview contains a valuable disclaimer saying that the results have not
been peer reviewed and should not be taken as established without independent
verification.  That disclaimer should remain.

The section describing the theoretical layer should be softened.  The current
language says that the theoretical layer establishes decidability and that the
two theoretical papers establish satisfiability iff existence of a valid finite
rank-\(d\) quotient.  Given the remaining obligations above, a safer statement is:
\begin{quote}
The theoretical layer proposes a route to decidability by reducing
satisfiability to the existence of a valid finite generated split-forest
quotient.  The current v2 manuscripts address the earlier density/Hasse and raw
pair-table defects, but the completeness extraction still depends on explicit
proofs of typed-\(\EQ\) separation, EQ-aware reachability, finite mosaic
locality, and coherent eventuality realization.
\end{quote}

\section{Suggested priority order for revision}

The following order would make the next revision maximally efficient.

\begin{enumerate}[leftmargin=2em]
\item \textbf{Fix equality data first.}  Separate \(\blk\) from \(\eqcong\),
  remove any implication from profile equality or reusable colour equality to
  semantic \(\EQ\), and make typed equality an explicit finite relation in the
  certificate.
\item \textbf{Make vertical reachability equality-class aware.}  Add mate-based
  existential and universal reachability clauses for \(\PP/\PPI\).
\item \textbf{Normalize \(\EQ\)-modalities.}  Either eliminate them by a
  preprocessing lemma or add a direct validity case.
\item \textbf{Prove finite mosaic locality.}  Replace the informal
  \(|\Cl(C_0)|+4\) claim by a theorem with an effective bound.
\item \textbf{Add the side-context lemma for \(\DR/\PO\) witnesses.}
\item \textbf{Prove simultaneous eventuality realization.}  This may be a direct
  lasso/branch proof rather than a Buechi automaton construction.
\item \textbf{Repair the TeX source.}  Define or remove \verb|\sst| and ensure
  clean compilation without nonstop-mode errors.
\end{enumerate}

\section{Conclusion}

The v2 updates fix the most visible structural problem in the earlier proof: the
proof no longer depends on a Hasse diagram or ambient root of an arbitrary
semantic \(\PP/\PPI\)-order.  The witness-generated submodel lemma and generated
cover reinterpretation put the argument much closer to the intended model class.
The support-closed mosaic idea also repairs the earlier raw pair-table weakness
in the right conceptual direction.

The decidability proof is therefore significantly more plausible than before.
However, it is not yet complete as a formal proof.  The key missing piece is a
fully specified finite certificate in which blocking, split-copy equality,
vertical reachability, and sibling mosaics interact coherently.  Once the
remaining obligations listed above are discharged, the original split-forest
strategy may well yield the desired decidability theorem.

\end{document}
